\section{Universality of the complete candidate subspace}
\label{sec:kernel-universality}

The critical nullity law alone does not determine the optimal test.
Two candidate subspaces of the same dimension can contain different
numbers of nondegenerate forms.  We therefore compare the laws of the
entire candidate subspaces, including their positions in the vector
space of alternating forms.

Throughout this section, $q$ is fixed, $n$ tends to infinity through
even integers,
\[
 E=\Alt(\F^n),\qquad D=\binom n2,\qquad t=D+b,
\]
and $b\in\mathbb Z$ is fixed.  The global offset notation is
$b=a-1$.  Under the null model, $K_1,\ldots,K_t$ are independent
uniform two-dimensional subspaces of $\F^n$, and
\[
 N=\{G\in E:G|_{K_i\times K_i}=0\text{ for every }i\}.
\]
Let $N^{\mathrm u}$ be the common kernel of $t$ independent uniform
linear functionals on $E$.  All total variation distances below
compare probability measures on the full lattice of subspaces of
the same $D$-dimensional space $E$.

\begin{theorem}[Complete kernel universality]
\label{ker:theorem}
For every fixed prime power $q$ and integer $b$,
\[
 \TV\bigl(\mathcal L(N),\mathcal L(N^{\mathrm u})\bigr)
 \longrightarrow0.
\]
Consequently, conditional on any fixed candidate dimension having
positive limiting probability, $N$ is asymptotically uniform on the
corresponding Grassmannian.  The conclusion holds in characteristic
two as well.
\end{theorem}

The proof has two parts.  We first compare, in aggregate, the
probabilities that each fixed-dimensional subspace survives all
observations.  A quantitative truncation of subspace M\"obius
inversion then upgrades those comparisons to total variation.  This
avoids an unproved assertion that a collection of moments uniquely
determines the limiting distribution.

\subsection{Exact subspace survival probabilities}

Write
\[
 \delta_n=(q+1)q^{-n}-q^{1-2n}.
\]
This is the probability that two independent uniform vectors
$u,v\in\F^n$ are linearly dependent.  Fix a nontrivial additive
character $\psi$ of $\F$.  Averaging first over $v$ gives
\begin{equation}
 \E_{u,v}\psi(G(u,v))=q^{-\operatorname{rank}G}
 \qquad(G\in E).
 \label{ker:rank-fourier}
\end{equation}
If $M\le E$ has dimension $j$, character orthogonality therefore
gives, before conditioning on independence,
\[
 \Pbb\{G(u,v)=0\text{ for all }G\in M\}
 =q^{-j}\sum_{G\in M}q^{-\operatorname{rank}G}.
\]
Every alternating form vanishes on a dependent pair.  Conditioning
on independence, which makes $\langle u,v\rangle$ a uniform
two-plane, proves the exact formula
\begin{equation}
 h(M):=\Pbb\{G|_{K\times K}=0\text{ for all }G\in M\}
 =\frac{q^{-j}\displaystyle\sum_{G\in M}
                 q^{-\operatorname{rank}G}-\delta_n}
        {1-\delta_n}.
 \label{ker:survival}
\end{equation}
In particular,
\begin{equation}
 \Pbb(M\le N)=h(M)^t,
 \qquad
 \Pbb(M\le N^{\mathrm u})=q^{-jt}.
 \label{ker:inclusion}
\end{equation}

We use the Gaussian binomial notation
\[
 \qbinom{m}{j}
 =\prod_{i=0}^{j-1}\frac{q^m-q^i}{q^j-q^i},
 \qquad
 g_j=\prod_{i=1}^j(1-q^{-i}),\qquad
 g_\infty=\prod_{i=1}^{\infty}(1-q^{-i}).
\]
As usual, the Gaussian binomial is zero outside $0\le j\le m$.
The empty product is one.

\begin{lemma}[Aggregate inclusion comparison]
\label{ker:aggregate}
For each fixed nonnegative integer $j$,
\begin{equation}
 \Delta_{n,j}:=
 \sum_{\substack{M\le E\\\dim M=j}}
 \left|h(M)^t-q^{-jt}\right|
 \longrightarrow0.
 \label{ker:aggregate-limit}
\end{equation}
For $j\ge1$, $t\ge1$, and any integer $0\le R<n$, define
\begin{align*}
 B_{n,R}&=\#\{G\in E:\operatorname{rank}G\le R\},\\
 \eta_q&=q^{-1}+(1-q^{-1})q^{-2}<1,\\
 A_{n,j,R}&=\frac{\eta_q+q^{j-R}}{1-\delta_n},\\
 \theta_{n,j,R}&=
 \frac{(q^j-1)(q^{-R}+\delta_n)}{1-\delta_n},\\
 m_{n,j}&=\qbinom Dj q^{-jt}.
\end{align*}
Then the following finite bound holds:
\begin{equation}
 \begin{split}
 \Delta_{n,j}\le{}&
 m_{n,j}\bigl(e^{t\theta_{n,j,R}}-1\bigr)\\
 &+A_{n,j,R}^{t}
       \sum_{s=1}^j B_{n,R}^{s}q^{-(j-s)b}
   +\sum_{s=1}^jB_{n,R}^{s}q^{-sD-jb}.
 \end{split}
 \label{ker:finite-aggregate}
\end{equation}
Moreover,
\begin{equation}
 B_{n,R}\le
 g_\infty^{-1}q^{D-\binom{n-R}{2}}
 \le g_\infty^{-1}q^{nR},
 \qquad
 m_{n,j}\le g_j^{-1}q^{-j(j+b)}.
 \label{ker:counts}
\end{equation}
\end{lemma}

\begin{proof}
The case $j=0$ is exact.  For the rank count, the radical of a form
of rank at most $R$ contains an $(n-R)$-space.  There are at most
$g_\infty^{-1}q^{R(n-R)}$ such spaces, and forms vanishing on one
descend to at most $q^{R(R-1)/2}$ alternating forms on its quotient.
This proves the first bound in \eqref{ker:counts}; its second
inequality follows by expanding the exponent.  The bound on
$m_{n,j}$ follows directly from the product formula for the Gaussian
binomial.

For a $j$-space $M$, let $s(M)$ be the dimension of the span of
\[
 \{G\in M:\operatorname{rank}G\le R\}.
\]
If $s(M)=0$, every nonzero element of $M$ has rank greater than $R$.
Equation \eqref{ker:survival} then gives
\[
 \left|\frac{h(M)}{q^{-j}}-1\right|
 \le\theta_{n,j,R}.
\]
For a nonnegative number $x$ with $|x-1|\le\theta$ and a positive
integer $t$, $|x^t-1|\le e^{t\theta}-1$.  Summing over these generic
$M$ gives the first term of \eqref{ker:finite-aggregate}.

Suppose next that $s(M)=s\ge1$.  Its low-rank span has a basis
consisting of low-rank forms.  Choose an ordered such basis, then
complete it by $j-s$ arbitrary vectors of $E$.  This overcounts the
possible $M$, and gives the upper bound
\begin{equation}
 \#\{M:\dim M=j,\ s(M)=s\}
 \le B_{n,R}^{s}q^{(j-s)D}.
 \label{ker:bad-count}
\end{equation}
Every nonzero alternating form has rank at least two.  All the
low-rank elements of $M$ lie in its specified $s$-space, so
\begin{align*}
 \sum_{G\in M}q^{-\operatorname{rank}G}
 &\le1+(q^s-1)q^{-2}+(q^j-q^s)q^{-R}\\
 &=q^s\bigl(q^{-s}+(1-q^{-s})q^{-2}\bigr)
      +(q^j-q^s)q^{-R}.
\end{align*}
The expression in parentheses is at most $\eta_q$.  Dropping the
negative term in the numerator of \eqref{ker:survival} therefore
shows that
\[
 h(M)\le q^{s-j}A_{n,j,R}.
\]
Together with \eqref{ker:bad-count}, this bounds the sum of
$h(M)^t$ over exceptional spaces by the second term of
\eqref{ker:finite-aggregate}.  Their $q^{-jt}$ contribution is bounded
by the third term.

Finally choose $\sigma$ with $\eta_q<\sigma<1$, put
$\beta=-\log_q\sigma>0$, and set
\[
 R=\left\lfloor\frac{\beta n}{8j}\right\rfloor.
\]
For all sufficiently large $n$, $A_{n,j,R}\le\sigma$ and
$t\theta_{n,j,R}\to0$.  For $1\le s\le j$, the exponent in the
second term satisfies
\[
 snR-\beta t
 \le\frac{\beta n^2}{8}
      -\beta\left(\binom n2+b\right),
\]
which tends to $-\infty$ quadratically.  The third term also tends
to zero quadratically.  This proves \eqref{ker:aggregate-limit}.
For fixed $q,j,b$, the first term is
$O_{q,j,b}(n^2q^{-\beta n/(8j)})$.
\end{proof}

The classification by the \emph{span} of low-rank elements is useful
here.  A small probability that any low-rank form survives does not,
by itself, control high moments of the number of surviving subspaces.
The aggregate estimate does provide that control.

\subsection{A quantitative inversion on the subspace lattice}

For a probability measure $\lambda$ on subspaces of $E$, put
$p_\lambda(M)=\Pbb_\lambda(M\le N)$.  For an integer $K\ge1$, define
the signed measure
\begin{equation}
 I_K(\lambda)(S)=
 \sum_{\substack{M\ge S\\\dim M\le K}}
 (-1)^{\dim M-\dim S}
 q^{\binom{\dim M-\dim S}{2}}p_\lambda(M)
 \quad(\dim S\le K),
 \label{ker:truncated-inversion}
\end{equation}
and set it to zero when $\dim S>K$.  Write $\|\cdot\|_1$ for the
sum of absolute masses, so that total variation between probabilities
is half this norm.

\begin{lemma}[Truncation bound]
\label{ker:inversion-bound}
Let
\[
 C_K=\prod_{i=1}^K(1+q^i).
\]
For every probability measure $\lambda$ on subspaces,
\begin{equation}
 \|\lambda-I_K(\lambda)\|_1
 \le C_K\E_\lambda\qbinom{\dim N}{K}
       +\Pbb_\lambda(\dim N>K).
 \label{ker:truncation-bound}
\end{equation}
For two such measures $\lambda,\lambda'$,
\begin{equation}
 \begin{split}
 \|I_K(\lambda)-I_K(\lambda')\|_1
 \le\sum_{j=0}^K c_j
 \sum_{\substack{M\le E\\\dim M=j}}
 |p_\lambda(M)-p_{\lambda'}(M)|,
 \end{split}
 \label{ker:inversion-comparison}
\end{equation}
where $c_j=\prod_{i=0}^{j-1}(1+q^i)$ and $c_0=1$.
\end{lemma}

\begin{proof}
For integers $r>0$ and $0\le L<r$, the finite $q$-binomial identity
\begin{equation}
 \sum_{j=0}^L(-1)^j q^{\binom j2}\qbinom rj
 =(-1)^Lq^{L(L+1)/2}\qbinom{r-1}{L}
 \label{ker:q-partial}
\end{equation}
follows by induction from the Gaussian-binomial recurrence.
For $L\ge r>0$, the corresponding full sum is zero; for $r=0$ it
is one.

Apply this identity to a deterministic $N$ of dimension $h$.  If
$h\le K$, \eqref{ker:truncated-inversion} reconstructs its point mass
exactly.  If $h>K$, the absolute mass at a subspace $S\le N$ of
dimension $s\le K$ is
\[
 q^{(K-s)(K-s+1)/2}\qbinom{h-s-1}{K-s}.
\]
Using
\[
 \qbinom hs\qbinom{h-s-1}{K-s}
 \le\qbinom hK\qbinom Ks
\]
and summing over $s$ bounds its total mass by
\[
 \qbinom hK
 \sum_{s=0}^K\qbinom Ks q^{(K-s)(K-s+1)/2}
 =\qbinom hK C_K.
\]
The equality is the finite $q$-binomial theorem.  The original point
mass, which lies above dimension $K$, contributes one more to the
error.  Averaging over $N$ proves \eqref{ker:truncation-bound}.

For \eqref{ker:inversion-comparison}, apply the triangle inequality
to \eqref{ker:truncated-inversion} and then interchange its sums.
A fixed $j$-space $M$ contributes the coefficient
\[
 \sum_{s=0}^j\qbinom js q^{\binom{j-s}{2}}
 =\prod_{i=0}^{j-1}(1+q^i)=c_j,
\]
again by the $q$-binomial theorem.
\end{proof}

\begin{proof}[Proof of Theorem~\ref{ker:theorem}]
Let $\mu_n=\mathcal L(N)$ and $\nu_n=\mathcal L(N^{\mathrm u})$.
Counting $K$-subspaces inside the random kernel gives
\[
 \E_{\nu_n}\qbinom{\dim N}{K}
 =\qbinom DK q^{-Kt}
 \longrightarrow
 m_K:=\frac{q^{-K(K+b)}}{g_K}.
\]
Lemma~\ref{ker:aggregate} implies that the corresponding expectation
under $\mu_n$ has the same limit.  Also,
\[
 \Pbb_\lambda(\dim N>K)
 \le\E_\lambda\qbinom{\dim N}{K}
\]
for either measure.  Lemma~\ref{ker:inversion-bound} and
\eqref{ker:aggregate-limit} now give, for each fixed $K\ge1$,
\begin{equation}
 \limsup_{n\to\infty}\|\mu_n-\nu_n\|_1
 \le2(C_K+1)m_K.
 \label{ker:limsup}
\end{equation}
Since $g_K\ge g_\infty>0$ and
\[
 C_K=q^{K(K+1)/2}\prod_{i=1}^K(1+q^{-i}),
\]
we have
\[
 C_Km_K
 =O_q\bigl(q^{-K(K-1)/2-Kb}\bigr)\longrightarrow0.
\]
Let $K\to\infty$ in \eqref{ker:limsup}.  This proves total variation
convergence.  The conditional assertion follows because $N^{\mathrm u}$
is exactly uniform among subspaces of a specified dimension and
conditioning on events with limiting probability bounded away from
zero preserves total variation convergence.
\end{proof}

\begin{corollary}[A finite computable comparison]
\label{ker:finite-TV}
For every $1\le K\le D$ and $t=D+b\ge1$,
\begin{equation}
 \begin{split}
 \TV(\mu_n,\nu_n)\le{}&
 \frac12\sum_{j=0}^Kc_j\Delta_{n,j}
 +(C_K+1)m_{n,K}
 +\frac{C_K+1}{2}\Delta_{n,K}.
 \end{split}
 \label{ker:finite-TV-bound}
\end{equation}
Here $m_{n,K}=\qbinom DK q^{-Kt}$, and the other quantities are
defined in Lemmas~\ref{ker:aggregate} and~\ref{ker:inversion-bound}.
Replacing every $\Delta_{n,j}$ by \eqref{ker:finite-aggregate}, with
an independently chosen integer $R_j\in[0,n-1]$ and the count bound
\eqref{ker:counts}, gives an entirely explicit upper bound.  It may
always be capped at one.  This is a certified computable comparison;
no numerical optimality of its constants is asserted.
\end{corollary}

\begin{proof}
The two Gaussian-binomial moments of order $K$ differ by at most
$\Delta_{n,K}$, and the uniform-row moment is exactly $m_{n,K}$.
Apply \eqref{ker:truncation-bound} to both measures, bound each tail
probability by its order-$K$ moment, and use
\eqref{ker:inversion-comparison}.  Dividing the resulting $\ell^1$
bound by two proves \eqref{ker:finite-TV-bound}.
\end{proof}

\subsection{Transfer of the planted experiment}

Let $Z(S)$ be the number of nondegenerate projective alternating forms
in $S$, and let $M_n$ denote the number of such forms in $E$.
If $H$ is any fixed nondegenerate form, the fraction of two-planes
isotropic for $H$ is
\begin{equation}
 p_n=\frac{q^{n-2}-1}{q^{n-1}-1}.
 \label{ker:isotropic-probability}
\end{equation}
For example, after choosing $u\ne0$, there are $q^{n-1}-q$
independent choices of $v$ with $H(u,v)=0$, out of $q^n-q$
independent choices in total.

In the planted experiment, a uniformly chosen nondegenerate
projective $H$ is shared by all observations, each of which is
uniformly $H$-isotropic.  Its likelihood ratio against the null is
exactly
\begin{equation}
 L_n(N)=\frac{Z(N)}{M_np_n^t}.
 \label{ker:actual-likelihood}
\end{equation}
Thus the entire candidate subspace is a sufficient statistic for
testing and for Bayesian recovery of the form.  Indeed, the
likelihood for each individual $H$ is $p_n^{-t}$ when $H\in N$ and
zero otherwise.  The posterior on the compatible nondegenerate
projective forms is therefore uniform.

There is a corresponding planted uniform-row experiment: choose
the same prior on $H$, then sample uniform linear functionals
annihilating $H$.  Its kernel likelihood against $\nu_n$ is
\begin{equation}
 L_n^{\mathrm u}(N)=\frac{Z(N)}{M_nq^{-t}}.
 \label{ker:uniform-likelihood}
\end{equation}

\begin{proposition}[Universality of both statistical laws]
\label{ker:planted-transfer}
The total variation distance between the two planted kernel laws
also tends to zero.  More strongly,
\begin{equation}
 \sum_{S\le E}Z(S)|\mu_n(S)-\nu_n(S)|\longrightarrow0.
 \label{ker:weighted-TV}
\end{equation}
Consequently the full null/planted pair of candidate-space laws is
asymptotically equivalent to the uniform-linear-constraint pair.
\end{proposition}

\begin{proof}
Put $h=\dim N$.  The elementary formulas
\[
 \qbinom h1=\frac{q^h-1}{q-1},\qquad
 \qbinom h2=\frac{(q^h-1)(q^h-q)}{(q^2-1)(q^2-q)}
\]
and Lemma~\ref{ker:aggregate}, with $j=1,2$, imply
\[
 \sup_n\bigl(\E_{\mu_n}q^{2h}+\E_{\nu_n}q^{2h}\bigr)<\infty
\]
after omission of finitely many initial $n$.  Since
$Z(S)\le(q^{\dim S}-1)/(q-1)$, the contribution to
\eqref{ker:weighted-TV} from $h>K$ is $O_{q,b}(q^{-K})$ uniformly
in $n$.  On $h\le K$ it is at most
$q^K\|\mu_n-\nu_n\|_1/(q-1)$, which tends to zero by
Theorem~\ref{ker:theorem}.  First taking $n\to\infty$, then
$K\to\infty$, proves \eqref{ker:weighted-TV}.

Finally $(qp_n)^t\to1$ because $t=O(n^2)$ and
$qp_n=1+O_q(q^{-n})$.  The elementary count of nondegenerate alternating
forms gives
\[
 M_nq^{-t}=\frac{q^{-b}}{q-1}
     \prod_{i=1}^{n/2}(1-q^{-(2i-1)}),
\]
which has a finite strictly positive limit at the fixed offset $b$.
Comparing the two densities
\eqref{ker:actual-likelihood}--\eqref{ker:uniform-likelihood}, using
\eqref{ker:weighted-TV} and the convergent normalizations, proves
the planted assertion.
\end{proof}

